% !PRINTED 617
\weberpage{617}
%% [page furniture, not body text -- source-confirmed 2026-08-10]
%% \S\,168. Reduction of the class equation in the cyclotomic fields.

%<<CONT
to determine. The first factor is the class number of the quadratic forms of the order $O'$ [\S\,100, (3)], which we denote by $H'$. We have therefore
\weq{3}{h = H' (\bar{E} O', \bar{E} A) = H' (O', A)}
(by \S\,100, 13.).

It remains to determine $(O', A)$.

If $\varrho_1, \varrho_2, \dots, \varrho_\mu$ denotes a complete system of representatives of $O'$ with respect to $A$, then the numbers
\weq{4}{N(\varrho_1),\ N(\varrho_2)\dots,\ N(\varrho_\mu)}
% [sic] printed N(\varrho_2)\dots, with no comma between N(\varrho_2) and the dots, a punctuation misprint in an otherwise comma-separated list (cf. \varrho_1, \varrho_2, \dots, \varrho_\mu in the preceding sentence)
must all be different with respect to the modulus $m$, and every norm of a number in $O$ is congruent with respect to $m$ to one of these numbers. If therefore $Z$ is the group of rational numbers, $M$ the group of those numbers $z$ from $Z$ which satisfy the condition
\[
z \equiv 1 \pmod{m}
\]
and $R_m$ the group of norm residues with respect to $m$, then (4) is a complete system of representatives of $R_m$ with respect to $M$, and consequently
\[
\mu = (O', A) = (R_m, M).
\]

For this one may also put (\S\,100, 12.):
\weq{5}{(O', A) = \frac{(Z, M)}{(Z, R_m)}.}

Now we have
\weq{6}{(Z, M) = \varphi(m) = m \Pi\Bigl(1 - \frac{1}{r}\Bigr),}
where $r$ runs through the prime factors of $m$, and consequently
\weq{7}{h = \frac{H'\varphi(m)}{(Z, R_m)}.}

We obtain the class field $\mathfrak{K}(A)$ of our group by adjoining to the class field $\mathfrak{K}(D)$ a primitive $m$th root of unity:
\weq{8}{\mathfrak{K}(A) = \mathfrak{K}(D, \varrho).}
For the condition that a prime ideal $\mathfrak{P}$ of the field $\mathfrak{K}(D,\varrho)$ dividing a prime $p$ be of the first degree is, if $(k)$ denotes a class invariant of the discriminant $D$:
\weq{9}{(k)^p \equiv (k), \quad \varrho^p \equiv \varrho \pmod{\mathfrak{P}},}
provided that all disturbing primes, e.g. those in the discriminants of the equations defining the numbers $(k)$ and $(\varrho)$, have been taken up into the excludent $S$.
